% chktex-file 8
%   Allow dashes
% chktex-file 17
%   Don't require matching number of brackets and parentheses
% chktex-file 11
%   Allow ellipsis
% chktex-file 44
%   Allow regex
% chktex-file 9
%   Allow unmatched brackets and parentheses
% chktex-file 36
%   ALlow parentheses without spaces in front of them
% chktex-file 40
%   Allow factorials
% chktex-file 12
%   Don't require inter-word spaces
% chktex-file 3
%   Don't require enclosing parentheses with braces
% chktex-file 13
%   Don't require intersentence spaces
% chktex-file 26
%   Allow spaces in front of punctuation
\documentclass{article}
\usepackage{amsmath}
\usepackage{amssymb}
\usepackage[a4paper, margin=1in]{geometry}
\usepackage{graphicx}
\usepackage{hyperref}
\usepackage{listings}
\usepackage{cancel}
\lstset{
basicstyle=\small\ttfamily,
columns=flexible,
breaklines=true
}

\begin{document}

\section*{Item Types Design}
\begin{itemize}
    \item We distinguish between `materials', `resources', and `equipment' items.
    \begin{itemize}
        \item \textbf{Materials}, like metals or chemicals, have a notion of quality and technological advancement, where certain techniques or
            facilities might produce/require certain tiers and material types. Many of the 4-6 tiers of each type are likely to require special
            resources found on exoplanets, along with new technologies. These types, in addition to the tier-less resources, should
            encompass all upkeep and the requirements of all production and construction.
        \item \textbf{Resources}, like oxygen or consumer goods, have some simple identity, and the only advancement is in terms of the scale or ease at which they are produced.
        \item \textbf{Equipment}, like a spacesuit, vehicle, deployable module, or building, is an individual item or construct with capabilities based on
            stats and type, and not handled in bulk/stacks. Some are formally `items' and can be carried (e.g. a spacesuit), while
            others are more abstract (e.g. a power plant or type of engineered microbe) and can use similar stat-like representations
            but are not inventory items.
    \end{itemize}
    \item The 6 following tiers are considered for materials. Only 1-3 are available at the start of the game.
    \begin{enumerate}
        \item \textbf{Baseline}: iron, copper, wood, stone capable of basic construction, especially ground-based construction within Earth-like conditions - used in modern housing, tools, power infrastructure, etc.
        \item \textbf{Modern}: steel, concrete, glass, silicon, etc. for more advanced construction, for bulk construction exposed to harsh conditions - used in standard modern factories, skyscrapers, tanks, digital infrastructure, etc.
        \item \textbf{High-end}: titanium, aluminum, carbon fiber, gold, rare earth metals, etc. for resilient ground and space-capable construction - used in high-end modern `smart' buildings, aircraft, spaceships, high-performance computing, etc.
        \item \textbf{Post-modern}: nanomaterials, aerogels, superconductors, etc. - modern research grade but not feasible at scale, likely allowing fusion, quantum computing, etc. if they became widely available
        \item \textbf{Future materials}: likely only theoretical or unconsidered at this point.
        \item \textbf{Future materials}: likely only theoretical or unconsidered at this point.
    \end{enumerate}
    \item The 7 following types are considered for materials.
    \begin{itemize}
        \item \textbf{Structural metals}: Used for structural support and machinery
        \begin{enumerate}
            \item Iron and other basic widespread metals that can be used for heavy construction
            \item Steel, magnesium, aluminum, etc. used in modern cities
            \item Titanium and advanced alloys used in `smart' construction and modern aerospace
            \item Post-modern - materials only at research stage now
            \item Future materials - likely only theoretical or unconsidered at this point.
            \item Future materials - likely only theoretical or unconsidered at this point.
        \end{enumerate}
        \item \textbf{Construction materials}: Non-metals used as structural base 
        \begin{enumerate}
            \item Wood, stone, etc. for basic widespread construction in earth-like conditions
            \item Concrete, glass, ceramics, etc. used at industrial scale
            \item Carbon fiber, aerogel, etc. non-metallic materials used in advanced construction and aerospace
            \item Post-modern - materials only at research stage now
            \item Future materials - likely only theoretical or unconsidered at this point.
            \item Future materials - likely only theoretical or unconsidered at this point.
        \end{enumerate}
        \item \textbf{Advanced materials}: Used for specialized properties like unusual conductivity, corrosion or heat resistance, etc.
        \begin{enumerate}
            \item Copper and other basic conductive materials for power infrastructure
            \item Silicon and other materials of modern computing chips, bulk solar panels, etc.
            \item Gold, silver, and rare earth metals sought after for high-end electronics and industrial applications
            \item Post-modern - materials only at research stage now
            \item Future materials - likely only theoretical or unconsidered at this point.
            \item Future materials - likely only theoretical or unconsidered at this point.
        \end{enumerate}
        \item \textbf{Fuels}: Used for energy and vehicles
        \begin{enumerate}
            \item Coal, wood, and other simple combustibles
            \item Oil, natural gas, etc. capable of at least fueling cars or generators
            \item Hydrazine, jet fuel, rocket fuel, etc. capable of fueling aircraft and spaceships
            \item Post-modern - materials only at research stage now
            \item Future materials - likely only theoretical or unconsidered at this point.
            \item Future materials - likely only theoretical or unconsidered at this point.
        \end{enumerate}
        \item \textbf{Chemicals}: Process reagents and products for manufacturing, fertilizers, pesticides, etc. Includes a wide variety of
            fertilizer, solvents, catalysts, polymers, rubbers, coatings, etc., but not directly biological
        \begin{itemize}
            \item Similar 1-6 scale
        \end{itemize}
        \item \textbf{Biomaterials}: Nutrients, pharmaceuticals, medical supplies, cell culture media, enzymes, engineered tissues, etc., but not
            directly industrial unless in a biotech context
        \begin{itemize}
            \item Similar 1-6 scale
        \end{itemize}
        \item \textbf{Nuclear materials}: Fissile materials for nuclear power
        \begin{itemize}
            \item Similar 1-6 scale
        \end{itemize}
    \end{itemize}
    \item \textbf{Resources}: Oxygen, GHGs, inert gases, toxic gases, water, food, consumer goods, energy
    \begin{itemize}
        \item Energy is unique in that, while being consumed and produced like a resource, it is not stored or distributed as a physical item.
    \end{itemize}
\end{itemize}

\end{document}

% Sample figure
%\begin{figure}[!ht]
%    \centering
%    \includegraphics[width=350px]{screenshot_1.png}
%    \caption{Sample caption}
%\end{figure}

% Sample hyperlink
% \item \href{URL}{Label}

% Sample listing
% \begin{lstlisting}[language=Python, caption=Sample Python Code]
% print("Hello World")
% \end{lstlisting}
